\section{Three dimensions, level counts and scope}
\label{sec:completion}

\subsection{The inherited three-dimensional case}

The preceding internal-window argument needs $n-1\ge3$ and is not
applied when $n=3$. In that dimension, Theorem~2.1 and its complete
table in C421~\cite{c421} classify every integral periodic orbit of
$(x,y,z)\mapsto(y,z,yz+a-x)$, with exact native least periods
\[
 \{1,2,3,4,5,6,8,9,12\}.
\]
That theorem is an explicitly inherited computer-assisted result,
not a new computation-free conclusion of this paper.

For clarity, its table gives precisely the structure needed here.
For fixed $a$, the constant words satisfy $a=2r-r^2$ and form a
finite set. The two-periodic alternating words satisfy
\[
 (u-1)(v-1)=1-a.
\]
If $a\ne1$, finite divisor enumeration gives all integer pairs;
if $a=1$, the solutions lie on $u=1$ or $v=1$. The exact-period
and orientation exclusions in that table are linear inequalities
or removal of finitely many parameters. Every other parametric
word is affine in its integer parameter, with linear range
conditions, and the two sporadic words give finitely many points.
Taking every consecutive triple of these words preserves this
description. An unrestricted integer parameter splits into its
nonnegative and negative rays, and a finite exclusion splits a
ray into finitely many rays or points. Thus all these sets are
effectively of the form \eqref{eq:linear-set}, with the exact
least-period labels supplied by the table. This proves the
$n=3$ part of Theorem~\ref{thm:main} using exactly the stated
prior theorem.

The C421 proof reduces its residual to a finite exact certificate,
with independently rebuilt cycles, and retains all unbounded
families symbolically. Its source, actual certificate and review
are available with that manuscript. None was rerun for the present
paper. In particular, our hand proof for $n\ge4$ must not be used
to relabel the three-dimensional imported theorem as computation-free.

\subsection{Finite levels and ordinary cycle counts}

Appendix~\ref{app:height} reproduces a previously completed auxiliary
height argument, including its degenerate branches. For any
$n\ge3$ and $a,D\in\Z$, it gives
\begin{equation}
 \|P\|_\infty\le R(n,a,D)
       :=2n^2\bigl(|a|+|D|+n+2\bigr)
 \label{eq:level-bound}
\end{equation}
for every periodic point on $K=D$. This is a level-dependent
bound, not a common bound over all levels.

Enumerate all integer points of $K=D$ in the box
$[-R(n,a,D),R(n,a,D)]^n$. Put an arrow $P\to F(P)$ when both
ends are in this finite set. Its directed cycles are exactly the
periodic points on the level: every genuine periodic orbit is inside
the box by \eqref{eq:level-bound}, and every graph cycle is an
actual native orbit. The cycle length is its least period. This
gives an exact terminating counting rule using only finite
integer arithmetic. The graph is a proved algorithm, not a
reported execution over all inputs.

Alternatively, use the atlas membership tests on this finite set.
Always count distinct coordinate points, not their possibly repeated
semilinear parameter representations. If $N_r(a,D)$ is the number
of points of least period $r$, then
$c_r(a,D)=N_r(a,D)/r$ is the number of oriented native cycles.
Consequently
\begin{align}
 \#\Fix(F^m\mid K=D,\Z^n)&=\sum_{r\mid m}r\,c_r(a,D),\\
 \zeta_{a,D}(z)&=\prod_r(1-z^r)^{-c_r(a,D)}.
 \label{eq:level-zeta}
\end{align}
There are only finitely many nonzero factors. These are ordinary
source-dynamical counts with no multiplicity weighting, prime
relabeling or altered clock. A finite-count zeta on the entire
integer lattice is not asserted.

For example, when $n=4$ and $a\ne0$, the previously observed
zero-product family $P=(0,a,0,t)$, $t\in\Z$, follows the word
\[
 (0,a,0,t,a,0,a,a-t).
\]
Every cyclic window of length three contains a zero, so this is
an actual $F$-word with period dividing eight. Four steps send
its first coordinate from $0$ to $a$, excluding all proper
divisors of eight; its least period is eight for every $t$.
Its invariant is $t(t-a)$, illustrating why an unbounded channel
can meet each fixed level in only finitely many points. This
family is an illustration, not the basis of the exhaustive atlas
or an additional contribution counted separately.

\subsection{Limitations and reproducibility}

The result concerns the specified single native map over the
ordinary integer lattice. It does not classify rational or complex
periodic schemes, arbitrary words in all Vieta involutions, or
unrelated additive forcing in each coordinate. The output is not
a canonical or minimal channel decomposition. The elementary
constant $L_n$ is large, and no useful complexity bound is claimed.

All new $n\ge4$ arguments, the classical finite constructions,
and the level-height proof are included here. The research-stage
bounded exploratory program is not a proof dependency; it neither
generated the full atlas nor replaces any infinite-parameter
argument. The inherited $n=3$ computation is identified above.
Source ownership and actual access limits are recorded in the
accompanying source audit. These bounded checks do not certify
worldwide priority or journal acceptance.

The ordinary factors in \eqref{eq:level-zeta} do not establish
target arithmetic Euler factors, root numbers, automorphy, a
spectral zero correspondence or a Hilbert--P\'olya realization.
The result is a source-system classification.

\paragraph{Availability and preparation.}
The LaTeX source, proof notes, source audit and internal review
records accompany this manuscript; the cited C421 package contains
its own exact computational evidence. Preparation and internal
checking used AI assistance; no human peer review is claimed.
No human participants or sensitive personal data are involved.
Named-author contribution, funding and competing-interest
information was not supplied, and no such declaration is inferred
from the anonymous working-paper format.
